\chapter{Results and Evaluation}
\label{chap:results}

% ============================================================================
% CHAPTER PURPOSE -- VERY HIGH PRIORITY
% Present evidence generated by Chapter \ref{chap:validation-method}, interpret
% it and answer the research questions. Separate observed result from
% explanation. Do not repeat methodology at length.
% ============================================================================

\section{Implementation and Verification Results}
\label{sec:verification-results}

\subsection{Software Test Results}
% Concise summary table: test family | count | result | notable issue.
% Detailed test cases belong in Appendix \ref{app:software-tests}.
% Report failures and corrections honestly; a defect found and fixed is useful
% verification evidence.

\subsection{Hand-Calculation Comparison}
% Present software versus independent calculation and numerical difference.
% Explain any rounding/indexing convention.

\subsection{Expected-Behaviour Results}
% Table/plots for limiting cases from Section
% \ref{sec:model-behaviour-verification}. Interpret whether the direction and
% magnitude are consistent with the equations and strategy reasoning.

\section{Historic Case Results}
\label{sec:historic-results}

% Use a consistent repeated structure for each case rather than free-form prose.
% For each:
% 1. race context and why selected;
% 2. decision state and input table;
% 3. predicted gain/post stop outcome;
% 4. actual observed result;
% 5. error and explanation;
% 6. annotated race trace.

\subsection{Successful Undercut Case(s)}
% HIGH PRIORITY. Include at least one clear success, but do not rely only on
% successes.

\subsection{Unsuccessful or Neutral Case(s)}
% HIGH PRIORITY. Demonstrating a negative prediction/outcome materially improves
% credibility and allows failure conditions to be discussed.

\subsection{Compromised or Boundary Case(s)}
% MEDIUM/HIGH PRIORITY. Traffic, warm-up, low degradation or unusual pace.
% Useful for showing limits of the minimum model and value of extensions.

\subsection{Cross-Case Validation Summary}
% MUST-HAVE TABLE:
% Case | predicted gain | observed gain | signed/absolute error | direction
% correct? | key confounder.
% Report aggregate values cautiously due to likely small sample size.

\section{Sensitivity Analysis Results}
\label{sec:sensitivity-results}

% VERY HIGH PRIORITY. Use a baseline case and one clear graph per parameter or a
% well-designed combined summary. Every graph requires:
% - baseline and range;
% - trend/threshold;
% - engineering explanation;
% - implication for strategy/model use.

\subsection{Tyre Degradation and Tyre Age}
% Examine both rate and current age where distinguishable.
% Identify threshold at which predicted gain changes sign if present.

\subsection{Compound Performance Delta}
% Explain interaction with offset length and warm-up.

\subsection{Offset-Lap Count / Target Stop Timing}
% This is likely a dominant mechanism; show cumulative nature clearly.

\subsection{Warm-Up Loss}
% Compare no warm-up, baseline and adverse warm-up assumptions.

\subsection{Relative Driver/Car Pace}
% HIGH PRIORITY to prevent strategy gain being mistaken for underlying pace.

\subsection{Traffic / Rejoin Penalty}
% Include if the term is retained in the final model.

\subsection{Sensitivity Ranking and Interactions}
% Summarise normalised/representative impact in a tornado/ranked table.
% MEDIUM PRIORITY: show one or two meaningful parameter interactions, such as
% degradation x warm-up or offset laps x compound delta. Do not create an
% exhaustive design if time is limited.

\section{Research Question Evaluation}
\label{sec:rq-evaluation}

% Answer each research question explicitly, with evidence citations to the
% relevant table/figure. Distinguish "supported for selected cases" from broad
% universal claims.

\section{Engineering Discussion}
\label{sec:engineering-discussion}

\subsection{Conditions Favouring an Undercut}
% Synthesise results: fresh-tyre delta, target degradation, offset duration,
% warm-up, clear rejoin and relative pace. Avoid simply restating graph slopes.

\subsection{Conditions Favouring Staying Out or an Overcut}
% Use negative/weak undercut results to explain low degradation, adverse warm-up,
% traffic and strong old-tyre pace. This demonstrates system understanding even
% if overcut is not fully implemented.

\subsection{Practical Interpretation of Model Output}
% Explain how an engineer should use G, post stop gap and race trace annotation:
% as decision support subject to assumptions, not an autonomous instruction.
% Discuss required margin when inputs are uncertain.

\subsection{Comparison with Existing Research}
% Return to Chapter \ref{chap:background}: where do observed trends agree or
% disagree with literature/accepted strategy mechanisms? Explain differences.

\section{Model Validity, Limitations and Generalisability}
\label{sec:validity-limitations}

% Organise rather than repeat a generic limitation list:
% - internal validity: implementation and parameter fitting;
% - construct validity: whether G captures "undercut effectiveness";
% - external validity: circuits/races/conditions beyond tested cases;
% - data limitations;
% - omitted dynamics and uncertainty.
% State which findings are robust and which are conditional.

\section{Objective Achievement Summary}
\label{sec:objective-achievement-results}

% Concise table: objective | evidence | status | caveat.
% This creates a direct bridge to Chapter \ref{chap:conclusion}.

\section{Chapter Summary}
\label{sec:results-summary}

% 1--2 paragraphs containing only the principal evidence-based findings.
% Do not introduce future work here; reserve it for Chapter 7.
